\documentclass[../DoAn.tex]{subfiles}
\begin{document}

Chương 1 đã trình bày bối cảnh và bài toán. Chương này tóm tắt ba khối kiến thức nền tảng cần để hiểu phần thiết kế sau: (1)~mô hình ngôn ngữ lớn dùng cho phân tích văn bản, (2)~các cách tiếp cận bài toán phân tích CV, và (3)~kỹ thuật so khớp CV--JD bằng embedding và RAG. Cuối chương là khảo sát các giải pháp hiện có và khoảng trống mà khóa luận hướng tới.

\section{Mô hình ngôn ngữ lớn và prompt}
\label{sec:llm}

Mô hình ngôn ngữ lớn (Large Language Model -- LLM) là loại mô hình được huấn luyện trên lượng văn bản rất lớn, có khả năng \textit{đọc một đoạn văn bản rồi tách ra các thông tin có cấu trúc} -- chính là điều khóa luận sử dụng để phân tích CV. Các mô hình tiêu biểu gồm ChatGPT (OpenAI) \cite{openai2023gpt4}, Gemini (Google), Claude (Anthropic); bên trong đều dựa trên kiến trúc \textit{Transformer} \cite{vaswani2017attention}. GPT-3 \cite{brown2020gpt3} là cột mốc chứng minh LLM có thể xử lý nhiều tác vụ khác nhau chỉ bằng vài ví dụ trong prompt (few-shot learning). Kiến trúc encoder BERT \cite{devlin2019bert} đặt nền tảng cho các mô hình embedding ngữ nghĩa, trong đó Sentence-BERT \cite{reimers2019sbert} mở rộng ra bài toán tính độ tương đồng câu -- tiền đề cho embedding JD/CV mà khóa luận sử dụng.

Để LLM thực hiện một tác vụ cụ thể, người dùng cần viết một câu hướng dẫn gọi là \textit{prompt}. Khóa luận thiết kế prompt theo hướng dẫn của \cite{white2023prompt}: gán vai trò cho mô hình, yêu cầu trả về dữ liệu JSON có cấu trúc cố định, và đưa ví dụ mẫu. Mô hình được dùng là \texttt{gpt-4o-mini} -- một phiên bản nhỏ và rẻ, đủ tốt cho phân tích CV.

\section{Phân tích CV bằng máy tính}

Phân tích CV là việc đọc một tệp CV (PDF, DOCX) và tách ra các trường có cấu trúc (họ tên, email, kinh nghiệm, học vấn, kỹ năng). Bài toán này khó vì: mỗi người viết CV theo một kiểu khác nhau, ranh giới giữa các phần đôi khi không rõ, tiếng Việt và tiếng Anh có quy ước khác nhau, và trích văn bản từ PDF/DOCX thường bị nhiễu (chữ rời, ngắt dòng bất thường).

Có ba cách tiếp cận phổ biến:

\textbf{Cách 1 -- Quy tắc (rule-based):} Dùng biểu thức chính quy (regex) để tìm các trường có định dạng cố định như email, số điện thoại, URL LinkedIn. Nhanh, miễn phí, ổn định, nhưng chỉ làm được với khuôn mẫu rõ.

\textbf{Cách 2 -- Học máy:} Huấn luyện mô hình nhận dạng tên riêng (công ty, chức vụ) trong văn bản \cite{chen2016resume,jiechieu2021skills}. Cần dữ liệu gán nhãn lớn, khó mở rộng sang ngôn ngữ mới.

\textbf{Cách 3 -- LLM:} Đưa cả CV cho LLM kèm prompt yêu cầu trích thông tin. Linh hoạt, hỗ trợ đa ngôn ngữ, nhưng tốn chi phí gọi API và chậm hơn.

Khóa luận chọn cách \textit{kết hợp} (hybrid): dùng regex cho các trường đơn giản (liên hệ, kỹ năng) và dùng LLM cho phần phức tạp (kinh nghiệm làm việc) -- tận dụng ưu điểm của cả hai và giảm chi phí API. Đánh giá định lượng của lựa chọn này được trình bày ở Chương~5.

\section{So khớp CV--JD: TF-IDF, Embedding và RAG}

Bài toán thứ hai của khóa luận -- tìm tin tuyển dụng (JD) phù hợp với CV của ứng viên từ một kho lớn -- thuộc nhóm \textit{Information Retrieval} \cite{manning2008ir}. Có hai họ phương pháp đo độ tương đồng giữa văn bản:

\textbf{TF-IDF (từ khoá):} Mỗi văn bản được biểu diễn bằng tần suất từ có chiết khấu nghịch đảo tần suất tài liệu \cite{sparck1972statistical}; độ tương đồng tính bằng cosine giữa hai vector. Nhanh, không tốn chi phí, nhưng không nắm được ngữ nghĩa -- ví dụ ``backend developer'' và ``software engineer'' bị coi là khác nhau dù gần nhau về ý. Biến thể BM25 \cite{robertson2009bm25} cải thiện công thức trọng số nhưng vẫn giữ bản chất dựa trên từ khoá.

\textbf{Embedding (ngữ nghĩa):} Mỗi văn bản được mã hoá thành một vector ngữ nghĩa dày đặc \cite{karpukhin2020dense} -- \texttt{text-embedding-3-small} của OpenAI \cite{openai2024embedding} tạo ra vector 1.536 chiều. Hai văn bản gần nghĩa nhau sẽ có vector gần nhau (cosine similarity cao \cite{turney2010frequency}) ngay cả khi từ khoá khác nhau. Đổi lại, cần gọi API để tính embedding và tốn không gian lưu trữ vector. Các hệ thống so khớp CV--JD \cite{qin2018enhancing,shen2018joint,he2021towards} xác nhận embedding vượt TF-IDF cho bài toán này khi từ ngữ giữa CV và JD không khớp trực tiếp.

\textbf{RAG (Retrieval-Augmented Generation):} Quy trình hai bước do Lewis et al. \cite{lewis2020rag} đề xuất -- bước \textit{retrieval} dùng phương pháp rẻ (TF-IDF/Embedding) lọc nhanh từ kho dữ liệu xuống vài chục kết quả, rồi mới dùng LLM (đắt nhưng tinh) để xếp hạng lại. Cách này tránh chi phí gọi LLM cho từng cặp (CV, JD) trong kho. Khóa luận áp dụng RAG hai bước cho bài toán tìm JD khớp yêu cầu; thiết kế chi tiết ở Chương~3.

\section{Các giải pháp hiện có và khoảng trống nghiên cứu}

Trên thị trường có nhiều công cụ liên quan đến CV nhưng đều thiếu một số yếu tố quan trọng. Bảng~\ref{tab:comparison} tóm tắt theo nhóm.

\begin{table}[H]
\centering
\caption{So sánh các giải pháp xây dựng CV hiện có}
\label{tab:comparison}
\begin{tabular}{|p{4.0cm}|p{4.0cm}|p{5.5cm}|}
\hline
\textbf{Nhóm} & \textbf{Đại diện} & \textbf{Hạn chế} \\ \hline
Nền tảng tuyển dụng Việt Nam & TopCV, VietnamWorks, JobsGO & Mạnh về CV Builder tiếng Việt; bước đầu tích hợp AI (TopCV CV Scoring 08/2025) nhưng không công bố phương pháp matching; ràng buộc apply trong hệ sinh thái của nền tảng \\ \hline
Công cụ tạo CV quốc tế & Canva, Zety, NovoResume, Enhancv & Chỉ cung cấp mẫu thiết kế hoặc gợi ý từ khoá; không phân tích sâu nội dung; ít hỗ trợ tiếng Việt \\ \hline
Dịch vụ phân tích CV thương mại & Sovren, Textkernel, DaXtra & Chi phí cao; phục vụ doanh nghiệp tuyển dụng; ít hỗ trợ tiếng Việt \\ \hline
\end{tabular}
\end{table}

\textbf{Khoảng trống nghiên cứu mà khóa luận hướng tới}: chưa có hệ thống mã nguồn mở phục vụ ứng viên Việt Nam mà công khai \textit{phương pháp} phân tích CV và so khớp CV--JD: các công cụ thương mại không công bố thuật toán matching và ràng buộc ứng viên ứng tuyển trong hệ sinh thái của nền tảng. Khóa luận trình bày cụ thể (i) thiết kế Hybrid Parser kết hợp regex và LLM, và (ii) quy trình RAG hai bước (Hybrid Retrieval + Structural/LLM Ranking) -- hai thiết kế có thể tái sử dụng cho các bài toán phân tích CV song ngữ Việt--Anh khác.

\end{document}
